\section{Significance and Broader Implications}
\label{sec:analysis-significance}

The preceding sections answer the research questions for the evaluated
configuration and bound what those answers can support. This section steps back
to reflect on two questions that the measurements raise but do not themselves
settle: how far the \emph{method} developed here transfers to other model
families, and what the decomposition it studies signifies beyond raw latency.
Both are framed as reasoned reflection bounded by the evidence rather than as
new empirical claims; their purpose is to make explicit the forward-looking
significance that the per-stage results otherwise leave implicit.

\subsection{Generalisability of the method to other neural architectures}
\label{sec:analysis-generalisability}

The headline numbers in this thesis are specific to ResNet-18
(\cref{sec:conclusion-limitations}), but the \emph{staged method} --- screen
split candidates under controlled local execution, explain them through
activation-transfer size and compute distribution
(\cref{sec:analysis-rq13}), then carry the selected decomposition through
local Kubernetes, managed cloud, and security-hardened deployment while changing
one layer at a time --- is largely architecture-agnostic. What determines
whether it transfers is not the milliseconds reported here but two structural
inputs it consumes at each stage: the set of \emph{candidate boundaries} and the
\emph{per-segment cost profile} (activation-transfer size and compute share) used
to screen them. It is therefore more useful to reflect on how those two inputs
change for a different architecture than to ask whether the absolute numbers
recur.

For another convolutional network the reasoning carries over directly. The
candidate boundaries remain the architecturally meaningful layer or block edges,
and the two-dimensional screening rule --- reject boundaries that expand the
activation budget the interface must carry, and reject boundaries that leave a
compute-degenerate downstream segment --- is unchanged. Only the cost profile
differs: a network without ResNet-18's channel-doubling, resolution-halving
structure would place its tolerable mid-region elsewhere, but the method
discovers that region empirically rather than assuming it, so no part of the
pipeline needs redesign.

The transformer case is the more substantive reflection, both because it was
raised explicitly during supervision and because the method's two inputs map
onto transformers cleanly despite the different architecture. The natural
candidate boundaries become the inter-block edges between encoder or decoder
layers, and at finer granularity the attention-versus-feed-forward sub-block
edges within a block, rather than residual stages. The activation-transfer
dimension still applies, but its driver changes: the tensor crossing a boundary
is the per-token hidden state, whose size scales with sequence length times
hidden dimension, so transfer cost becomes sequence-length-dependent in a way it
is not for a fixed-resolution image model --- a boundary that is inexpensive for
a short prompt can dominate for a long context. The compute-distribution
dimension also still applies, but its shape inverts: transformer blocks are far
more uniform than ResNet stages, with approximately equal parameter and
floating-point cost per block, so the compute-degeneracy failure mode that
rejected \texttt{split\_after\_layer4} in \cref{sec:analysis-rq11} is far less
likely to bind and a near-balanced split is achievable at almost any block
edge. The reflective consequence is that, run on a transformer, the same method
would likely report a flatter suitability landscape in which the boundary choice
is governed mainly by activation transfer --- and hence by sequence length ---
rather than by compute balance, the opposite weighting to the late-CNN case
studied here. The screening, explanatory, deployment, and hardening stages
themselves would execute unchanged; only the cost profile they consume differs.

Three aspects would require genuine extension rather than reinterpretation, and
naming them identifies the concrete modifications the pipeline would need. First,
the fixed single-input-shape benchmark contract (\cref{sec:methods-contract})
would have to be parameterised over sequence length, because a single shape no
longer characterises the workload. Second, autoregressive decoding introduces
key-value cache state that persists across tokens, whereas the stateless
per-request chain measured here carries no such cross-step state; a split placed
inside an attention block would have to account for where that cache lives and
whether it crosses the boundary. Third, secure-transformer inference is a
distinct research line in which the protection mechanism is cryptographic
(for example homomorphic or multi-party protocols) rather than the
infrastructure-level identity, transport, and VM-level controls measured in this
thesis; that changes the cost model at the boundary entirely and is not a simple
substitution into the present pipeline. These are extensions of the method, not
merely caveats on it: each identifies a specific place where the staged design
would have to be enriched to remain valid for a transformer workload.

\subsection{Significance of per-layer service decomposition}
\label{sec:analysis-impact}

This thesis frames its contribution defensively, as a bounded overhead budget for
the architectural, deployment, and security properties that a monolithic baseline
does not provide (\cref{sec:analysis-evaluation}). That framing is deliberate and
evidence-bounded, but it understates a capability that the decomposition opens and
that the monolith structurally cannot offer, and it is worth stating that
significance plainly. When each architectural segment runs as an independently
addressable service behind a typed inter-layer interface, the service boundary
becomes a point of control.

The most immediate consequence is \emph{per-layer policy}. Each boundary is a
place at which identity, authorisation, encryption, confidential execution, audit,
and placement can be enforced independently of the rest of the model. This thesis
already demonstrates the principle: communication-level identity, mutual TLS, and
authorisation policy were applied at the inter-service path
(\cref{sec:analysis-rq21}), and VM-level confidential execution was applied
\emph{selectively} to the single downstream segment that receives intermediate
activations rather than to the whole model (\cref{sec:analysis-rq22}). A monolith
admits none of this granularity: it is one trust domain, one placement decision,
and one policy. Decomposition turns ``the model'' into a set of independently
governable units, and the overhead this thesis measures is precisely the price of
that governance capability rather than a cost incurred for its own sake.

A more speculative future direction is \emph{cross-organisational
composition}. Once layers are independently addressable services joined by an
explicit interface, nothing in the architecture requires them to be owned or
operated by the same organisation. Different segments could be provided by
different parties that agree to interoperate, so that a model is assembled at run
time from components hosted under distinct administrative and legal domains ---
analogous to a ``digital embassy'', in which one party's workload runs within
another's jurisdiction under agreed terms. The per-layer controls above are what
make such a scenario credible rather than merely conceivable: a participating
organisation can require an authenticated identity from its upstream, run its own
segment inside a confidential VM so that the (possibly foreign) host cannot
observe the activations it processes, and admit traffic only from authorised
peers. The thesis does not implement this scenario, but it is a plausible extension
of the decomposition the thesis measures, and it reframes the act of splitting a
monolith into services as one that \emph{enables} new governance, jurisdiction,
and selective-protection options rather than one that only adds overhead.

Both consequences depend on the boundary between layers being a stable, explicit
contract rather than an internal implementation detail, which is where the
thesis's measurements bear most directly on future design. The activation-transfer
interface used here --- a typed tensor message with an explicit shape and a
functional-parity contract, carried over gRPC --- is exactly such a boundary, and
the central explanatory finding that boundary cost is jointly determined by
activation-transfer size and compute placement (\cref{sec:analysis-rq13}) is
informative for how an inter-layer API for composed inference would have to be
parameterised. Such an API would need to expose at least the activation tensor's
shape and dtype, because those determine both the transfer cost and the size of
the confidentiality envelope that must protect the activation in transit and at
rest, together with enough of the compute-distribution contract for a consumer to
reason about where its segment sits in the chain and what it is responsible for
executing. The thesis does not design this API, but its results indicate which
properties such an interface would have to treat as first-class.

These reflections should not be read as claims that the thesis has demonstrated
per-layer governance or cross-organisational inference at scale; it has not. What
it has demonstrated is that the decomposition on which those capabilities depend
carries a bounded and characterised overhead across controlled, orchestrated,
cloud, and security-hardened settings. That bounded answer is the precondition
for the larger questions: decomposition is worth pursuing as an enabler of
governance and composition precisely because its cost has been shown to be
contained rather than prohibitive, which is the sense in which ``measuring the
overhead'' is foundational rather than incidental to the broader significance of
the work.
